% Template: Đường tròn tâm O bán kính R với dây cung AB
% Dùng cho: hình học đường tròn, góc nội tiếp, tiếp tuyến
% --------------------------------------------------------
\documentclass[border=5pt]{standalone}
\usepackage[utf8]{inputenc}
\usepackage[T5]{fontenc}
\usepackage[vietnamese]{babel}
\usepackage{tikz}
\usepackage{helvet}
\renewcommand{\familydefault}{\sfdefault}
% ── Hệ màu tập trung (chỉnh tại đây để đổi theme) ──────────
\usepackage{xcolor}
\definecolor{mauchinh} {HTML}{1565C0}
\definecolor{mauphu}   {HTML}{43A047}
\definecolor{maunoibat}{HTML}{E53935}
\definecolor{mauvien}  {HTML}{424242}
{mauchinh} {HTML}{1565C0}
{mauphu}   {HTML}{43A047}
{maunoibat}{HTML}{E53935}
{mauvien}  {HTML}{424242}
\usetikzlibrary{calc, angles, quotes, arrows.meta}

\begin{document}
\begin{tikzpicture}[
    scale=1.3,
    >=Stealth,
    line join=round,
    line cap=round,
    every node/.style={font=\sffamily\small},
    thick,
]
    % Bán kính R = 2
    \pgfmathsetmacro{\R}{2}
    % Góc của điểm A trên đường tròn
    \pgfmathsetmacro{\angA}{130}
    % Góc của điểm B trên đường tròn
    \pgfmathsetmacro{\angB}{30}

    % Tâm đường tròn O
    \path (0, 0) coordinate (O);
    % Điểm A trên đường tròn tại góc 130°
    \path (\angA:\R) coordinate (A);
    % Điểm B trên đường tròn tại góc 30°
    \path (\angB:\R) coordinate (B);

    % Đường tròn tâm O bán kính R
    \draw (O) circle (\R);

    % Dây cung AB
    \draw[mauchinh] (A) -- (B);

    % Bán kính OA (nét đứt)
    \draw[dashed, thin] (O) -- (A);
    % Bán kính OB (nét đứt)
    \draw[dashed, thin] (O) -- (B);

    % Nhãn tâm O
    \path (O) node[below left] {$O$};
    % Nhãn điểm A
    \path (A) node[above left] {$A$};
    % Nhãn điểm B
    \path (B) node[right] {$B$};

    % Điểm tại O, A, B
    \foreach \p in {O, A, B} {
        \fill (\p) circle (2pt);
    }
\end{tikzpicture}
\end{document}
